\documentclass[10pt,a4paper]{amsart}
\usepackage[margin=19mm]{geometry}
\usepackage{amsmath,amssymb,mathtools}
\usepackage[scr=rsfs,cal=euler]{mathalfa}
\usepackage{xfrac}
\usepackage[unicode,hidelinks,bookmarks=false]{hyperref}
\newcommand{\ie}{i.e.\ }
\newcommand{\cf}{cf.\ }
\newcommand{\resp}{respectively\ }
\newcommand{\Subsection}{\subsection}
\providecommand{\nobreakdash}{\nobreak\hbox{-}}
\setlength{\emergencystretch}{2em}
\multlinegap0pt
\usepackage{bm}
\usepackage{bbm}
\usepackage{dsfont}
\usepackage{xspace}
\usepackage{cancel}
\usepackage{mathtools}
\usepackage{amsthm}
\makeatletter
\@ifundefined{lemma}{\newtheorem{lemma}{Lemma}}{}
\makeatother
\newcommand*{\defeq}{\mathrel{\vcenter{\baselineskip0.5ex \lineskiplimit0pt
                     \hbox{\scriptsize.}\hbox{\scriptsize.}}}%
                     =}
\newcommand{\bM}{\mathbf{M}}
\newcommand{\bu}{\mathbf{u}}
\newcommand{\bv}{\mathbf{v}}
\newcommand{\bw}{\mathbf{w}}
\newcommand{\bzeta}{\boldsymbol{\zeta}}
\newcommand{\der}{\mathrm{d}}
\newcommand{\vertiii}[1]{{\left\vert\kern-0.25ex\left\vert\kern-0.25ex\left\vert #1 
    \right\vert\kern-0.25ex\right\vert\kern-0.25ex\right\vert}}
\renewcommand{\d}{\, \mathrm{d}}
\renewcommand{\leq}{\leqslant}
\title[One-Dimensional Carrollian Fluids III: Global Existence and ]{One-Dimensional Carrollian Fluids III: Global Existence and Weak Continuity in $L^\infty$}
\author{P. Marios Petropoulos, Simon Schulz, Grigalius Taujanskas}
\date{Source: arXiv:2407.05972v1 (CC BY 4.0)}
\begin{document}
\maketitle

\noindent\emph{Source and licence.} One-Dimensional Carrollian Fluids III: Global Existence and Weak Continuity in $L^\infty$. Version arXiv:2407.05972v1 (CC BY 4.0), distributed under the Creative Commons Attribution 4.0 International licence, as recorded in the arXiv licence metadata for this exact version and shown on the abstract page https://arxiv.org/abs/2407.05972v1. Full rights notice: licenses/arxiv-2407.05972-CC-BY-4.0.txt.\par\smallskip

\noindent\textit{Editorial scope.} Counted unit: the Lemma whose source label is \texttt{lem:fixed pt i} in the article (source0.tex lines 1056--1058, source label \texttt{lem:fixed pt i}), delivered together with the article's own complete original proof of it (source0.tex lines 1060--1239, one \texttt{proof} environment of 180 lines). No mathematics is written, repaired or re-derived: statement and proof are the article's own text line for line. Required context retained uncounted: eq:op norm of auxiliary M (lines 950--952). Every definition, display and label of those lines is retained unchanged. Omitted from this package and not redistributed: 1--949, 953--1055, 1059--1059, 1240--1705. Nothing inside a retained excerpt is omitted or altered: every retained body is delivered exactly as the article's lines are written, so no omission receipt of any kind accompanies this package. Citations: the retained excerpts carry no citation command at all, so no bibliography entry of the article is transcribed and no reference is redistributed. Reference adaptation: none was needed and none was made; every reference inside the delivered unit resolves inside it, so no unresolved reference or citation is produced. Printed slips kept literal: no printed word, symbol, hypothesis or number of the counted statement or of its proof was altered. Layout and class: the article's own class is \texttt{amsart} with packages including graphicx, amsmath, amsfonts, amssymb, amsthm, tcolorbox, dsfont, xcolor, hyperref, cleveref, tabularx, float, nicefrac, biblatex; the wrapper is the lane's pinned \texttt{amsart} editorial wrapper at 10pt on A4, which restates the theorem-like environments the retained text needs and adds the article's own macro definitions that the retained text needs (\texttt{\textbackslash bM}, \texttt{\textbackslash bu}, \texttt{\textbackslash bv}, \texttt{\textbackslash bw}, \texttt{\textbackslash bzeta}, \texttt{\textbackslash d}, \texttt{\textbackslash defeq}, \texttt{\textbackslash der}, \texttt{\textbackslash leq}, \texttt{\textbackslash vertiii}), with extra wrapper packages (bm, bbm, dsfont, xspace, cancel, mathtools). The delivered numbering is the unit's own. Assets: no retained span contains a figure environment, a graphics-inclusion command, a PDF-page inclusion, a table or a TikZ picture; no image file is copied into this package, the package declares no asset dependency and no image file is redistributed with it. Licence: version arXiv:2407.05972v1 of this article is distributed under the Creative Commons Attribution 4.0 International licence, as recorded in the arXiv licence metadata for this exact version and shown on the abstract page https://arxiv.org/abs/2407.05972v1. A full rights notice is delivered at licenses/arxiv-2407.05972-CC-BY-4.0.txt. Characters: every retained excerpt is pure ASCII, so the delivered engine, XeLaTeX with its default Latin Modern text font, reports no missing character.
\begin{equation}\label{eq:op norm of auxiliary M}
    \vertiii{\bM(\bu)} \leq C \big( |\phi_1(\bu)| + |\phi_2(\bu)| \big) \leq C 
\end{equation}

\begin{lemma}\label{lem:fixed pt i}
    The mapping $L : X \to X$ is continuous and compact.
\end{lemma}

\begin{proof}
We divide the proof into several steps. 

\smallskip 


 \noindent 1. \textit{$L$ is self-mapping, i.e.~$L(X) \subset X$.}  Let $\bw \defeq L\bv$, and note that there holds 
\begin{equation*}
    \left\lbrace\begin{aligned}
    &\bw_t - \varepsilon\bw_{xx} = - \bM(\bv)\bv_x, \\ 
    &\bw|_{t=0} = \bu_0 
    \end{aligned}\right. 
\end{equation*}
in the weak sense. In turn, performing the standard parabolic estimate by testing the equation with $\bw$, we obtain 
\begin{equation*}
    \frac{1}{2}\frac{\der}{\der t}\int_\mathbb{R} |\bw|^2 \d x + \varepsilon \int_\mathbb{R} |\bw_x|^2 \d x = -\int_{\mathbb{R}} \bw \cdot \bM(\bv)\bv_x \d x \leq \frac{1}{2}\int_\mathbb{R} |\bw|^2 \d x + C\int_\mathbb{R}|\bv_x|^2 \d x, 
\end{equation*}
and Gr\"onwall's Lemma implies, for a.e.~$t \in (0,T)$, 
\begin{equation*}
    \Vert \bw (t,\cdot) \Vert^2_{L^2(\mathbb{R})} \leq \Vert \bu_0 \Vert^2_{L^2(\mathbb{R})} \exp\Big( C \int_0^t \Vert \bv_x(s,\cdot) \Vert^2_{L^2(\mathbb{R})} \d s \Big). 
\end{equation*}
It therefore follows that 
\begin{equation}\label{eq:L2H1 est fixed point}
    \begin{aligned} \Vert \bw \Vert^2_{L^\infty(0,T;L^2(\mathbb{R}))} + & \varepsilon \Vert \bw_x \Vert^2_{L^2((0,T)\times\mathbb{R})} \\ 
    &\leq C(1+T)\Vert \bu_0 \Vert^2_{L^2(\mathbb{R})} \exp\Big( C \Vert \bv_x \Vert^2_{L^2((0,T)\times\mathbb{R})}  \Big) + CT \Vert \bv_x \Vert^2_{L^2((0,T)\times\mathbb{R})}, 
    \end{aligned}
\end{equation}
whence $L$ is self-mapping since $L^\infty(0,T;L^2(\mathbb{R})) \subset L^2(0,T;L^2(\mathbb{R}))$ for $T$ finite. 


\smallskip 
\noindent 2. \textit{Estimate on time derivatives}. We begin with an estimate on the time derivative which will enable us to employ the Aubin--Lions Lemma. Let $\bzeta \in L^2(0,T;H^1(\mathbb{R}))$. Testing against the equation, we obtain 
\begin{equation*}
   \begin{aligned}
       |\langle \bw_t , \bzeta \rangle| &\leq \varepsilon \Big| \int_0^T \int_\mathbb{R} \bw_x \cdot \bzeta_x \d x \d t \Big| + \Big| \int_0^T \int_\mathbb{R} \bzeta \cdot \bM(\bv) \bv_x \d x \d t \Big| \\ 
       &\leq \varepsilon \Vert \bw_x \Vert_{L^2((0,T)\times\mathbb{R})} \Vert \bzeta_x \Vert_{L^2((0,T)\times\mathbb{R})} + C\Vert \bzeta \Vert_{L^2((0,T)\times\mathbb{R})} \Vert \bv_x \Vert_{L^2((0,T)\times\mathbb{R})}, 
   \end{aligned} 
\end{equation*}
whence 
\begin{equation*}
   \begin{aligned}
       \Vert \bw_t \Vert_{L^2(0,T;H^{-1}(\mathbb{R}))} \leq \varepsilon \Vert \bw_x \Vert_{L^2((0,T)\times\mathbb{R})} + C\Vert \bv_x \Vert_{L^2((0,T)\times\mathbb{R})}, 
   \end{aligned} 
\end{equation*}
\textit{i.e.}, using \eqref{eq:L2H1 est fixed point}, 
\begin{equation}\label{eq:time deriv est fixed pt}
   \begin{aligned}
       \Vert \bw_t \Vert_{L^2(0,T;H^{-1}(\mathbb{R}))} \leq &\sqrt{\varepsilon}C(1+T)^{\frac{1}{2}}\Vert \bu_0 \Vert_{L^2(\mathbb{R})} \exp\Big( C \Vert \bv_x \Vert^2_{L^2((0,T)\times\mathbb{R})} \big) \\ 
       &+ C(1+T^{\frac{1}{2}}) \Vert \bv_x \Vert_{L^2((0,T)\times\mathbb{R})}. 
   \end{aligned} 
\end{equation}

\smallskip 
\noindent 3. \textit{Weak compactness in $X$}. Let $\{\bv_n\}_n$ be a bounded sequence in $X$ and define $\bw_n  \defeq  L\bv_n$ for all $n$. The Aubin--Lions Lemma implies that, up to a subsequence which we do not relabel, there holds
\begin{equation}
    \bv_n \to \bv \quad \text{weakly in } X, \text{ strongly in } L^2((0,T)\times\mathbb{R}), \text{ and a.e.} 
\end{equation}
for some $\bv \in X$, where we also used Alaoglu's Theorem and the fact that $X$ is reflexive. Furthermore, the estimates \eqref{eq:L2H1 est fixed point} and \eqref{eq:time deriv est fixed pt} imply that $\{\bw_n\}_n$ also satisfies the assumptions of the Aubin--Lions Lemma, so there exists $\bw \in X$ such that, up to a subsequence which we do not relabel, using also Alaoglu's Theorem, 
\begin{equation}\label{eq:conv from aubin lions}
    \bw_n \to \bw \quad \text{weakly in } X, \text{ strongly in } L^2((0,T)\times\mathbb{R}), \text{ and a.e.} 
\end{equation}

\smallskip 
\noindent 4. \textit{Strong compactness in $L^2(0,T;H^1(\mathbb{R}))$}. It remains to show that the convergence also occurs strongly with respect to the topology of $X$, \textit{i.e.}~it remains to show that $\lim_n \Vert (\bw_n)_x - \bw_x \Vert_{L^2((0,T)\times\mathbb{R})} = 0$ and an analogous strong convergence for the time derivatives in the space $L^2(0,T;H^{-1}(\mathbb{R}))$. Concerning the first convergence, since we already know that $\{\bw_n\}_n$ converges weakly to $\bw$ in $L^2(0,T;H^1(\mathbb{R}))$, it suffices to show 
\begin{equation}\label{eq:required strong conv for cpct fixed pt}
    \lim_n \Vert (\bw_n)_x \Vert_{L^2((0,T)\times\mathbb{R})} = \Vert \bw_x \Vert_{L^2((0,T)\times\mathbb{R})}. 
\end{equation}
The weak formulation for $\bw_n$ reads 
\begin{equation}\label{eq:weak form wn fixed pt}
   -\int_0^T \int_\mathbb{R} \bw_n \cdot \bzeta_t \d x \d t + \varepsilon \int_0^T \int_\mathbb{R} (\bw_n)_x \cdot \bzeta_x \d x \d t = - \int_0^T \int_\mathbb{R} \bM(\bv_n) (\bv_n)_x \cdot \bzeta \d x \d t 
\end{equation}
for all $\bzeta \in C^\infty([0,T]\times\mathbb{R})$ such that $\bzeta(0,\cdot)=\bzeta(T,\cdot) = 0$. The matrix $\bM$ is continuous, so $\bv_n \to \bv$ a.e.~implies that $\bM(\bv_n) \to \bM(\bv)$ a.e. Furthermore, the bound \eqref{eq:op norm of auxiliary M} gives $\vertiii{\bM(\bv_n)} \leq C$, whence $C|\bzeta| \in L^1((0,T)\times\mathbb{R})$ is a suitable dominating function for $\bzeta^\intercal \bM(\bv_n)$ which converges a.e.~to $\bzeta^\intercal \bM(\bv)$. The Dominated Convergence Theorem then yields 
\begin{equation*}
    \bzeta^\intercal \bM(\bv_n) \to \bzeta^\intercal \bM(\bv) \quad \text{strongly in } L^p((0,T)\times\mathbb{R}) \text{ for all } p \in [1,\infty). 
\end{equation*}
Since $(\bv_n)_x \to \bv_x$ weakly in $L^2((0,T)\times\mathbb{R})$, it follows that there holds the convergence for the product 
\begin{equation*}
    \bM(\bv_n) (\bv_n)_x \cdot \bzeta \rightharpoonup \bM(\bv) \bv_x \cdot \bzeta \quad \text{weakly in } L^1((0,T)\times\mathbb{R}). 
\end{equation*}
Using the above and the aforementioned weak convergences, we deduce from \eqref{eq:weak form wn fixed pt} that there holds 
\begin{equation}\label{eq:weak form w fixed pt}
   -\int_0^T \int_\mathbb{R} \bw \cdot \bzeta_t \d x \d t + \varepsilon \int_0^T \int_\mathbb{R} \bw_x \cdot \bzeta_x \d x \d t = - \int_0^T \int_\mathbb{R} \bM(\bv) \bv_x \cdot \bzeta \d x \d t. 
\end{equation}
A standard approximation argument allows us to introduce $\bw_n$ into the weak formulation \eqref{eq:weak form wn fixed pt} and $\bw$ into \eqref{eq:weak form w fixed pt} on an arbitrary time interval $[t_1,t_2] \subsetneq [0,T]$. Subtracting the two resulting equations, we get 
\begin{equation}\label{eq:diff of weak forms fixed pt}
   \begin{aligned} \frac{1}{2}&\int_\mathbb{R} \big( |\bw_n(t_2,x)|^2 -  |\bw(t_2,x)|^2 \big) \d x - \frac{1}{2}\int_\mathbb{R} \big(|\bw_n(t_1,x)|^2 - |\bw(t_1,x)|^2 \big) \d x \\ 
   &= \varepsilon \int_{t_1}^{t_2} \int_\mathbb{R} \big( |(\bw_n)_x|^2 - |\bw_x|^2 \big) \d x \d t -\int_{t_1}^{t_2}\int_\mathbb{R} \Big( \bM(\bv_n) (\bv_n)_x \cdot \bv_n - \bM(\bv) \bv_x \! \cdot \! \bv \Big) \d x \d t. 
   \end{aligned}
\end{equation}
Recall that the convergence $\bw_n \to \bw$ strongly in $L^2((0,T)\times\mathbb{R})$ implies, up to a subsequence which we do not relabel, 
\begin{equation*}
    \bw_n(t,\cdot) \to \bw(t,\cdot) \quad \text{strongly in } L^2(\mathbb{R}) \text{ for a.e.~}t. 
\end{equation*}
It follows that the left-hand side of \eqref{eq:diff of weak forms fixed pt} vanishes in the limit as $n\to\infty$. Hence, using also the continuity of $s \mapsto |s|$ to interchange the limit with the absolute value, 
\begin{equation}\label{eq:diff of weak forms fixed pt checkpoint}
      \begin{aligned} \lim_n \Big| \Vert (\bw_n )_x \Vert^2_{L^2((t_1,t_2)\times\mathbb{R})} -&  \Vert \bw_x \Vert^2_{L^2((t_1,t_2)\times\mathbb{R})}\Big| \\ 
      &= \varepsilon^{-1} \bigg|\lim_n \int_{t_1}^{t_2}\int_\mathbb{R} \Big( \bM(\bv_n) (\bv_n)_x \cdot \bv_n - \bM(\bv) \bv_x \cdot \bv \Big) \d x \d t\bigg|. 
   \end{aligned}
\end{equation}
For the term on the right-hand side, we write
\begin{equation}\label{eq:using dct to conclude fixed pt cpctness}
    \begin{aligned}
      \Vert \bv_n^\intercal \bM(\bv_n) - \bv^\intercal \bM(\bv) \Vert_{L^2} &\leq \Vert \bM(\bv_n) \cdot (\bv_n - \bv) \Vert_{L^2} + \Vert \big( \bM(\bv_n) - \bM(\bv) \big) \cdot \bv \Vert_{L^2} \\ 
      &\leq C \Vert \bv_n - \bv \Vert_{L^2} + \bigg( \int_{0}^{T}\int_\mathbb{R} \vertiii{\bM(\bv_n) - \bM(\bv)}^2 |\bv|^2 \d x \d t \bigg)^{\frac{1}{2}}. 
    \end{aligned}
\end{equation}
The first term on the right-hand side manifestly vanishes due to the strong convergence in $L^2$ for the sequence $\{\bv_n\}_n$. Recall that $\bM(\bv_n) \to \bM(\bv)$ a.e.~ and $\vertiii{\bM(\bv_n) - \bM(\bv)}^2 \leq C$, so the Dominated Convergence Theorem implies that the second term on the right-hand side vanishes in the limit as $n\to\infty$. It follows that 
\begin{equation*}
    \begin{aligned}
      \bv_n^\intercal \bM(\bv_n) \to \bv^\intercal \bM(\bv) \quad \text{strongly in } L^2((0,T)\times\mathbb{R}). 
    \end{aligned}
\end{equation*}
Since $(\bv_n)_x \rightharpoonup \bv_x$ weakly in $L^2((0,T)\times\mathbb{R})$, it follows that 
\begin{equation*}
    \bM(\bv_n) (\bv_n)_x \cdot \bv_n \rightharpoonup \bM(\bv) \bv_x \cdot \bv \quad \text{weakly in } L^1((0,T)\times\mathbb{R}), 
\end{equation*}
and hence 
\begin{equation*}
      \begin{aligned} 
      \lim_n \int_{t_1}^{t_2}\int_\mathbb{R} \Big( \bM(\bv_n) (\bv_n)_x \cdot \bv_n - \bM(\bv) \bv_x \cdot \bv \Big) \d x \d t = 0. 
   \end{aligned}
\end{equation*}
Returning to \eqref{eq:diff of weak forms fixed pt checkpoint}, we obtain 
\begin{equation*}
    \lim_n \Vert (\bw_n)_x \Vert_{L^2((t_1,t_2)\times\mathbb{R})} = \Vert \bw_x \Vert^2_{L^2((t_1,t_2)\times\mathbb{R})} \quad \text{for a.e.~}t_1, \, t_2. 
\end{equation*}
We now relax the assumption $[t_1,t_2] \subsetneq [0,T]$ (by taking a further subsequence if necessary), and deduce 
\begin{equation*}
    \lim_n \Vert (\bw_n)_x \Vert_{L^2((0,T)\times\mathbb{R})} = \Vert \bw_x \Vert^2_{L^2((0,T)\times\mathbb{R})}; 
\end{equation*}
this, in addition to the weak convergence in $L^2(0,T;H^1(\mathbb{R}))$, implies 
\begin{equation}\label{eq:strong conv L2H1 for fixed pt}
    \bw_n \to \bw \quad \text{strongly in } L^2(0,T;H^1(\mathbb{R})). 
\end{equation}


\smallskip 
\noindent 5. \textit{Strong compactness of time derivatives in $L^2(0,T;H^{-1}(\mathbb{R}))$}. It remains to show that $(\bw_n)_t \to \bw_t$ strongly in $L^2(0,T;H^{-1}(\mathbb{R}))$. Taking the difference of the two weak formulations, there holds 
\begin{equation*}
    \begin{aligned}
        |\langle (\bw_n)_t \! & - \! \bw_t , \bzeta \rangle| \\ 
        &\leq \varepsilon \int_0^T \int_\mathbb{R} \big| \big( (\bw_n)_x \! -\! \bw_x \big) \cdot \bzeta_x \big| \d x \d t + \int_0^T \int_\mathbb{R} \big| \big(\bM(\bv_n)(\bv_n)_x \! -\! \bM(\bv)\bv_x \big) \! \cdot \! \bzeta \big| \d x \d t \\ 
        &\leq \Big(\varepsilon \Vert \bw_n \! - \! \bw \Vert_{L^2(0,T;H^1(\mathbb{R}))} \!+\! \Vert \bM(\bv_n)(\bv_n)_x \!-\! \bM(\bv)\bv_x
 \Vert_{L^2((0,T)\times\mathbb{R})}\Big)\Vert \bzeta \Vert_{L^2(0,T;H^1(\mathbb{R}))}
    \end{aligned}
\end{equation*}
for all $\bzeta \in L^2(0,T;H^1(\mathbb{R}))$. We compute 
\begin{equation*}
    \begin{aligned}
        \Vert \bM(\bv_n)(\bv_n)_x - \bM(\bv)&\bv_x
 \Vert_{L^2((0,T)\times\mathbb{R})} \\ &\leq \Vert \bM(\bv_n)(\bv_n - \bv)_x \Vert_{L^2((0,T)\times\mathbb{R})} + \Vert (\bM(\bv) - \bM(\bv_n))\bv_x
 \Vert_{L^2((0,T)\times\mathbb{R})} \\ 
 &\leq C \Vert \bv_n \! - \! \bv \Vert_{L^2(0,T;H^1(\mathbb{R}))} + \Vert (\bM(\bv_n)\! -\! \bM(\bv) ) \bv_x\Vert_{L^2((0,T)\times\mathbb{R})}
    \end{aligned}
\end{equation*}
so that there holds 
\begin{equation*}
    \begin{aligned}
        \Vert (\bw_n)_t - \bw_t \Vert_{L^2(0,T;H^{-1}(\mathbb{R}))} \leq \, & \varepsilon \Vert \bw_n - \bw \Vert_{L^2(0,T;H^1(\mathbb{R}))} + C\Vert \bv_n - \bv \Vert_{L^2(0,T;H^1(\mathbb{R}))} \\ 
        &+ \Vert (\bM(\bv_n)\! -\! \bM(\bv) ) \bv_x\Vert_{L^2((0,T)\times\mathbb{R})}, 
    \end{aligned}
\end{equation*}
giving
\begin{equation*}
    \limsup_n  \Vert (\bw_n)_t - \bw_t \Vert_{L^2(0,T;H^{-1}(\mathbb{R}))} = \limsup_n \Vert (\bM(\bv_n)\! -\! \bM(\bv) ) \bv_x\Vert_{L^2((0,T)\times\mathbb{R})}. 
\end{equation*}
Arguing as we did to control the second term on the right-hand side of the final line of \eqref{eq:using dct to conclude fixed pt cpctness}, we use the Dominated Convergence Theorem to deduce
\begin{align*}
    \limsup_n \Vert (\bM(\bv_n)\! -\! \bM(\bv) ) \bv_x\Vert_{L^2((0,T)\times\mathbb{R})} 
    &\leq \limsup_n \bigg(\int_0^T \int_\mathbb{R} \vertiii{\bM(\bv_n) - \bM(\bv)}^2 |\bv_x|^2 \d x \d t \bigg)^{\frac{1}{2}} \\ 
    &= 0, 
\end{align*}
and it follows that 
\begin{equation}\label{eq:time deriv cpctness fixed pt}
    (\bw_n)_t \to \bw_t  \quad \text{strongly in } L^2(0,T;H^{-1}(\mathbb{R})). 
\end{equation}

\smallskip 
\noindent 6. \textit{Compactness in $X$ and continuity}. By combining \eqref{eq:strong conv L2H1 for fixed pt} and \eqref{eq:time deriv cpctness fixed pt}, we have shown that, given any bounded sequence $\{\bv_n\}_n$ in $X$, the sequence $\{L\bv_n \}_n$ admits a strongly convergent subsequence in $X$, \textit{i.e.}~the mapping $L:X \to X$ is compact. The continuity of $L$ is immediate from the compactness and the uniqueness of limits, \textit{i.e.} given a sequence $\{\bv_n \}_n$ in $X$ converging to $\bv \in X$, the previous proof shows that $\{\bw_n=L\bv_n\}_n$ converges to $\bw=L\bv \in X$. 
\end{proof}

\end{document}
